% Project 99b R109 APA author-date reference list (author-review metadata only).
% Publisher/source metadata links; NOT the unpublished manuscript.
% Cite the primary papers themselves, not this listing.
\bibitem[Alghamdi and Alghizzi(2026)]{ref39} Alghamdi, L. H., \& Alghizzi, T. M. (2026). Conditional validity in LLM-mediated L2 assessment: An argument-based systematic review and meta-analysis. Frontiers in Research Metrics and Analytics, 11, 1908592. \url{https://doi.org/10.3389/frma.2026.1908592}
\bibitem[Anghel et al.(2026a)]{ref49} Anghel, C., Anghel, A. A., Cocu, A., Craciun, M. V., Moldovanu, S., Vlase, M., \& Sandu, C. (2026a). RubricAdapt-LLM: Measuring criterion-level adaptation and score shifts under alternative grading rubrics. Informatics, 13(9), 149. \url{https://doi.org/10.3390/informatics13090149}
\bibitem[Anghel et al.(2026b)]{ref50} Anghel, C., Anghel, A. A., Craciun, M. V., Cocu, A., Moldovanu, S., Sandu, C., Dragosloveanu, C. D. M., Andrei, C. A., Vulpe, D.-E., Maier, C., Scheau, C., Dragosloveanu, S., \& Potop, V. (2026b). GradeDrift-LLM: Measuring student-history-induced score drift in LLM-based automated grading. Machine Learning and Knowledge Extraction, 8(8), 252. \url{https://doi.org/10.3390/make8080252}
\bibitem[Anghel et al.(2026c)]{ref47} Anghel, C., Anghel, A. A., Craciun, M. V., Cocu, A., Vulpe, D.-E., Andrei, C. A., Maier, C., Scheau, C., Dragosloveanu, S., \& Cergan, R. (2026c). GradeAgentOps: A verification-first framework for evidence-anchored LLM exam grading. AI, 7(6), 198. \url{https://doi.org/10.3390/ai7060198}
\bibitem[Anghel et al.(2026d)]{ref48} Anghel, C., Anghel, A. A., Vlase, M., Craciun, M. V., Cocu, A., Andrei, C. A., Vulpe, D.-E., Dragosloveanu, S., Scheau, C., Maier, C., \& Potop, V. (2026d). ConsensusGrade: A human-variability-aware framework for evaluating LLM-based automated grading. Applied System Innovation, 9(9), 174. \url{https://doi.org/10.3390/asi9090174}
\bibitem[Anghel et al.(2026e)]{ref20} Anghel, C., Craciun, M. V., Cocu, A., Anghel, A. A., Balau, A. S., Istrate, A., \& Anghele, A.-D. (2026e). EvalHack: Answer-Side Prompt Injection for Probing LLM Exam-Grading Panel Stability. Information, 17(3), 297. \url{https://doi.org/10.3390/info17030297}
\bibitem[Anghel et al.(2026f)]{ref19} Anghel, C., Craciun, M. V., Cocu, A., Anghel, A. A., \& Istrate, A. (2026f). SteadyEval: Robust LLM Exam Graders via Adversarial Training and Distillation. Computers, 15(1), 55. \url{https://doi.org/10.3390/computers15010055}
\bibitem[Apillanes et al.(2026)]{ref61} Apillanes, J., Weber, J. L., Gago-Masague, S., Wong-Ma, J., \& Yeh, T. Y. (2026). Fighting Fire with Fire: LLM-Assisted Grading of Handwritten CS Assessments. In Proceedings of the 57th ACM Technical Symposium on Computer Science Education V.1 (SIGCSE TS 2026), pp. 66--72. \url{https://doi.org/10.1145/3770762.3772615}
\bibitem[Aytutuldu et al.(2025)]{ref27} Aytutuldu, I., Yol, O., \& Akgul, Y. S. (2025). Integrating LLMs for Grading and Appeal Resolution in Computer Science Education. arXiv preprint arXiv:2504.13557. \url{https://arxiv.org/abs/2504.13557}
\bibitem[Azad et al.(2020)]{ref57} Azad, S., Chen, B., Fowler, M., West, M., \& Zilles, C. (2020). Strategies for Deploying Unreliable AI Graders in High-Transparency High-Stakes Exams. In Artificial Intelligence in Education (AIED 2020), LNCS 12163, pp. 16--28. \url{https://doi.org/10.1007/978-3-030-52237-7_2}
\bibitem[Brevik et al.(2026)]{ref35} Brevik, A., Jerpseth, H., \& Lafontan, S. R. (2026). Large language models as grading assistants in public health education: A method-comparison study of essay-style exam assessment. Frontiers in Education, 11, 1904452. \url{https://doi.org/10.3389/feduc.2026.1904452}
\bibitem[Cai(2026)]{ref51} Cai, Y. (2026). Prompt injection attacks on educational large language models for higher and vocational education. Scientific Reports, 16, article 15594. \url{https://doi.org/10.1038/s41598-026-46563-1}
\bibitem[\c{C}all\i{} and Er(2026)]{ref26} \c{C}all\i{}, E., \& Er, E. (2026). When technological momentum overshadows pedagogical alignment: a systematic review of AI-generated formative feedback in higher education. Assessment \& Evaluation in Higher Education. \url{https://doi.org/10.1080/02602938.2026.2674230}
\bibitem[Chiang et al.(2024)]{ref16} Chiang, C.-H., Chen, W.-C., Kuan, C.-Y., Yang, C., \& Lee, H.-Y. (2024). Large language model as an assignment evaluator: Insights, feedback, and challenges in a 1000+ student course. EMNLP, 2489--2513. \url{https://doi.org/10.18653/v1/2024.emnlp-main.146}
\bibitem[Cope et al.(2026)]{ref59} Cope, K. L., Frankenreiter, J., Hirst, S., Posner, E. A., Schwarcz, D., \& Thorley, D. (2026). Grading Machines: Can AI Exam-Grading Replace Law Professors? Journal of Law and Empirical Analysis, 3(1), 2--22. \url{https://doi.org/10.1177/2755323X261434265}
\bibitem[Crompton and Burke(2026)]{ref03} Crompton, H., \& Burke, D. (2026). Higher education assessments in the age of GenAI: A systematic review. International Journal of Educational Technology in Higher Education. \url{https://doi.org/10.1186/s41239-026-00627-4}
\bibitem[Cruz et al.(2026)]{ref10} Cruz, L. F., da Silva, M. M., \& Mamede, H. S. (2026). AI-assisted, instructor-supervised grading and feedback in higher education: Design and evaluation of an end-to-end pipeline. Computers and Education Open. \url{https://doi.org/10.1016/j.caeo.2026.100411}
\bibitem[Deepshikha(2025)]{ref41} Deepshikha, D. (2025). A comprehensive review of AI-powered grading and tailored feedback in universities. Discover Artificial Intelligence, 5, article 251. \url{https://doi.org/10.1007/s44163-025-00517-0}
\bibitem[Denny et al.(2024)]{ref55} Denny, P., Smith, D. H., Fowler, M., Prather, J., Becker, B. A., \& Leinonen, J. (2024). Explaining Code with a Purpose: An Integrated Approach for Developing Code Comprehension and Prompting Skills. Proceedings of ITiCSE 2024, 283--289. \url{https://doi.org/10.1145/3649217.3653587}
\bibitem[D'Urso and Sciarrone(2026)]{ref24} D'Urso, S., \& Sciarrone, F. (2026). Criterium: Assisting Teachers in Fair and Consistent Grading of Open-Ended Questions. In Learning Technologies and Systems: ICWL/SETE 2025 Revised Selected Papers (LNCS 16425, pp. 83--95). Springer. \url{https://doi.org/10.1007/978-981-92-0042-9_7}
\bibitem[Flod\'{e}n(2025)]{ref06} Flod\'{e}n, J. (2025). Grading exams using large language models: A comparison between human and AI grading of exams in higher education using ChatGPT. British Educational Research Journal, 51, 201--224. \url{https://doi.org/10.1002/berj.4069}
\bibitem[Gao et al.(2024)]{ref38} Gao, R., Merzdorf, H., Anwar, S., Hipwell, M., \& Srinivasa, A. (2024). Automatic assessment of text-based responses in post-secondary education: A systematic review. Computers \& Education: Artificial Intelligence, 6, 100206. \url{https://doi.org/10.1016/j.caeai.2024.100206}
\bibitem[Goulas et al.(2026)]{ref45} Goulas, S., Megalokonomou, R., \& Sotirakopoulos, P. (2026). Why do experts miss AI's errors? Evidence from a randomized labeling experiment. PNAS Nexus, 5(6), pgag146. \url{https://doi.org/10.1093/pnasnexus/pgag146}
\bibitem[Gr\'evisse(2024)]{ref58} Gr\'evisse, C. (2024). LLM-based automatic short answer grading in undergraduate medical education. BMC Medical Education, 24, 1060. \url{https://doi.org/10.1186/s12909-024-06026-5}
\bibitem[Habibullah et al.(2026)]{ref36} Habibullah, A., Alshoibi, Y., Alshiekh, M., Khan, S., \& Khan, N. (2026). Where LLM graders succeed and break: Evidence from two computer-science exams. arXiv preprint arXiv:2609.29333. \url{https://arxiv.org/abs/2609.29333}
\bibitem[Heinrich et al.(2025)]{ref43} Heinrich, T., Baily, S., Chen, K.-W., DeOliveira, J., Park, S., \& Wang, N. C.-H. (2025). AI-assisted grading and personalized feedback in large political science classes: Results from randomized controlled trials. PLOS ONE, 20(8), e0328041. \url{https://doi.org/10.1371/journal.pone.0328041}
\bibitem[Humble(2026)]{ref22} Humble, N. (2026). Ethical Implications of Prompt Injection in AI-mediated Grading: An Adversarial Red-Team Evaluation. Journal of Academic Ethics, 24, article 85. \url{https://doi.org/10.1007/s10805-026-09758-9}
\bibitem[Jin et al.(2026)]{ref02} Jin, Z., Li, X., Chen, X., Zhang, H., \& Dai, L. (2026). AI-driven assessment in higher education: a systematic review. Interactive Learning Environments. \url{https://doi.org/10.1080/10494820.2026.2711402}
\bibitem[Jukiewicz and Wyrwa(2026)]{ref25} Jukiewicz, M., \& Wyrwa, M. (2026). Can ChatGPT Replace the Teacher in Assessment? A Review of Research on the Use of Large Language Models in Grading and Providing Feedback. Applied Sciences, 16(2), 680. \url{https://doi.org/10.3390/app16020680}
\bibitem[Kane(2013)]{ref28} Kane, M. T. (2013). Validating the Interpretations and Uses of Test Scores. Journal of Educational Measurement, 50(1), 1--73. \url{https://doi.org/10.1111/jedm.12000}
\bibitem[Kay et al.(2026)]{ref08} Kay, W. P., Rutherford, S. M., Smith, D. P., Shore, A. M., \& Francis, N. J. (2026). Can GenAI be trained to mimic human markers of extended written assignments in higher education? Assessment \& Evaluation in Higher Education. \url{https://doi.org/10.1080/02602938.2026.2710688}
\bibitem[Kerwat et al.(2026)]{ref46} Kerwat, L., Donaldson, D. L., \& Mahomed, A. (2026). Can large language models reproduce higher education grade bands? Cross-model study of calibration and grading bias in authentic student writing. Assessment \& Evaluation in Higher Education. \url{https://doi.org/10.1080/02602938.2026.2734796}
\bibitem[Klapproth(2026)]{ref29} Klapproth, F. (2026). Redesigning Teaching, Learning, and Assessment in the Age of Generative AI: A Constructive Alignment Perspective for Higher Education. Education Sciences, 16(9), 1507. \url{https://doi.org/10.3390/educsci16091507}
\bibitem[Korthals et al.(2026)]{ref21} Korthals, L., Akrong, E., Geller, G., Rosenbusch, H., Grasman, R., \& Visser, I. (2026). Towards Reliable LLM Grading Through Self-Consistency and Selective Human Review: Higher Accuracy, Less Work. Machine Learning and Knowledge Extraction, 8(3), 74. \url{https://doi.org/10.3390/make8030074}
\bibitem[Latif et al.(2024)]{ref23} Latif, A., ElBarachi, M., Tsaramirsis, G., \& Khan, F. Q. (2024). Generative AI enabled robotic process automation: A practical case study. In 2024 IEEE 30th International Conference on Parallel and Distributed Systems (ICPADS), pp. 769--775. \url{https://doi.org/10.1109/ICPADS63350.2024.00104}
\bibitem[Li, Filippov et al.(2026)]{ref37} Li, H., Filippov, F., Lin, Y., He, P., Yang, K., Chu, Y., Cui, Y., Liu, H., \& Tang, J. (2026). ``Important You Should Give Me Full Credits!'': Exploring prompt injection attacks on LLM-based automatic grading systems. Findings of the Association for Computational Linguistics: EMNLP 2026, 43057--43071. \url{https://aclanthology.org/2026.findings-emnlp.2024/}
\bibitem[Li, Kurimchak et al.(2026)]{ref33} Li, W., Kurimchak, M., Andres, A., \& Whitmer, J. (2026). Evaluating the landscape of automated scoring with GenAI: A systematic review. Proceedings of the Impactful and Responsible AI Systems for Education Workshop, PMLR, 339, 140--145. \url{https://proceedings.mlr.press/v339/li26a.html}
\bibitem[Lin et al.(2026)]{ref42} Lin, M. P.-C., Chang, D. H., Winne, P. H., Bowen, G. M., \& Ryoo, J. (2026). Mapping generative AI in higher education teaching and learning: A scoping review. International Journal of Educational Technology in Higher Education, 23, Article 55. \url{https://doi.org/10.1186/s41239-026-00630-9}
\bibitem[MDPI(2026)]{ref32} MDPI. (2026, accessed 1 November). Research and Publication Ethics: Artificial Intelligence policy. \url{https://www.mdpi.com/ethics}
\bibitem[Milani et al.(2025)]{ref17} Milani, A., Franzoni, V., Florindi, E., Omarbekova, A., Bekmanova, G., \& Yergesh, B. (2025). When AI Is Fooled: Hidden Risks in LLM-Assisted Grading. Education Sciences, 15(11), 1419. \url{https://doi.org/10.3390/educsci15111419}
\bibitem[Orazbek et al.(2026)]{ref01} Orazbek, Y., Shildibekov, Y., \& Balkenova, A. (2026). AI-based automated grading in higher education: A scoping review and landscape analysis for sustainable education. Procedia Computer Science, 290, 902--907. \url{https://doi.org/10.1016/j.procs.2026.09.005}
\bibitem[Page et al.(2021)]{ref04} Page, M. J., McKenzie, J. E., Bossuyt, P. M., Boutron, I., Hoffmann, T. C., Mulrow, C. D., Shamseer, L., Tetzlaff, J. M., Akl, E. A., Brennan, S. E., Chou, R., Glanville, J., Grimshaw, J. M., Hr\'objartsson, A., Lalu, M. M., Li, T., Loder, E. W., Mayo-Wilson, E., McDonald, S., \ldots{} Moher, D. (2021). The PRISMA 2020 statement: An updated guideline for reporting systematic reviews. BMJ, 372, n71. \url{https://doi.org/10.1136/bmj.n71}
\bibitem[Papageorgiou and Pierrakeas(2026)]{ref30} Papageorgiou, K., \& Pierrakeas, C. (2026). Do AI Grading Systems Systematically Differ from Human Teachers' Grading? Evidence of Bias and Consistency in Educational Assessment. Education Sciences, 16(7), 1147. \url{https://doi.org/10.3390/educsci16071147}
\bibitem[Pecuchova et al.(2025)]{ref05} Pecuchova, J., Benko, L., \& Drlik, M. (2025). Automated grading of open-ended questions in higher education using GenAI models. International Journal of Artificial Intelligence in Education, 35, 3813--3846. \url{https://doi.org/10.1007/s40593-025-00517-2}
\bibitem[Poli\v{c}ar et al.(2025)]{ref44} Poli\v{c}ar, P. G., \v{S}pendl, M., Curk, T., \& Zupan, B. (2025). Automated assignment grading with large language models: Insights from a bioinformatics course. Bioinformatics, 41(Suppl. 1), i21--i29. \url{https://doi.org/10.1093/bioinformatics/btaf196}
\bibitem[Saez et al.(2026)]{ref34} Saez, Y., Garcia, L. M., Mochon, A., \& Isasi, P. (2026). Evaluating large language models for AI-assisted grading: A framework and case study in higher education. Scientific Reports, 16, 18035. \url{https://doi.org/10.1038/s41598-026-48656-3}
\bibitem[Savage and Rebello(2025)]{ref13} Savage, S., \& Rebello, N. S. (2025). Using an LLM to investigate students' explanations on conceptual physics questions. Physics Education Research Conference Proceedings, 399--404. \url{https://doi.org/10.1119/perc.2025.pr.Savage}
\bibitem[Sebastian et al.(2025)]{ref14} Sebastian, S. K., Chakravarthi, R., Sreedharan, J., Mathew, L., Nosseir, N., Ebby, S., \& Ramakrishnan, S. (2025). Enhancing Medical Education through LLM: Integration of Generative GPT for Assessment and Feedback. Medical Science Educator, 35, 29--31. \url{https://doi.org/10.1007/s40670-024-02258-8}
\bibitem[Speiser(2025)]{ref56} Speiser, S. (2025). Assessing the Real-World Impact of Disagreement Between Human Graders and LLMs. In Proceedings of the 40th ACM/SIGAPP Symposium on Applied Computing (SAC 2025), pp. 48--53. \url{https://doi.org/10.1145/3672608.3707736}
\bibitem[Sun et al.(2026)]{ref60} Sun, W., Chen, L., Cai, Y., Xie, H., Zeng, Y., \& Zhang, Y. (2026). EDU-CIRCUIT-HW: Evaluating Multimodal Large Language Models on Real-World University-Level STEM Student Handwritten Solutions. Findings of ACL 2026, 15281--15314. \url{https://doi.org/10.18653/v1/2026.findings-acl.751}
\bibitem[Talmi et al.(2026)]{ref07} Talmi, D., Benn, Y., Tibon, R., Corsi, G., \& Marcoci, A. (2026). AI in University Assessment: Evaluating the Opportunities and Risks of Automated Marking. University of Cambridge. Report: \url{https://www.ai.cam.ac.uk/reports/ai-in-university-assessment-evaluating-the-opportunities-and-risks-of-automated-marking/} ; official sampling and prompt-selection note: \url{https://www.cam.ac.uk/stories/ai-university-essay-grading}
\bibitem[Ting and O'Briain(2026)]{ref53} Ting, Y.-S., \& O'Briain, T. (2026). Teaching Astronomy with Large Language Models. The Open Journal of Astrophysics, 9. \url{https://doi.org/10.33232/001c.160220}
\bibitem[University of Birmingham(n.d.)]{ref31} University of Birmingham. (n.d.). Principles on the Use of Artificial Intelligence Tools to Support Grading and the Provision of Student Feedback. University policy webpage, accessed 1 November 2026. \url{https://www.birmingham.ac.uk/libraries/education-excellence/gai/principles-on-the-use-of-ai}
\bibitem[Wanjura et al.(2026)]{ref18} Wanjura, B., Shapiro, D., Mollick, E., Mollick, L., \& Meincke, L. (2026). Prompting Science Report 5: This Is an Excellent Paper: The Effects of Prompt Injection on Grading. Wharton Generative AI Labs. \url{https://gail.wharton.upenn.edu/research-and-insights/hidden-prompt-injections/}
\bibitem[Wetzler et al.(2025)]{ref09} Wetzler, E. L., Cassidy, K. S., Jones, M. J., Frazier, C. R., Korbut, N. A., Sims, C. M., Bowen, S. S., \& Wood, M. (2025). Grading the Graders: Comparing Generative AI and Human Assessment in Essay Evaluation. Teaching of Psychology, 52(3), 298--304. \url{https://doi.org/10.1177/00986283241282696}
\bibitem[Wu et al.(2026)]{ref52} Wu, H.-N., Chu, M.-N., \& Hsu, J.-L. (2026). Comparing GPT and human raters in essay assessment: Variability, bias, and the potential of LLM-based scoring. Computers and Education Open, 10, 100341. \url{https://doi.org/10.1016/j.caeo.2026.100341}
\bibitem[Yang(2024)]{ref11} Yang, Y. (2024). The Reliability of Using ChatGPT in Rating EFL Writings. Shanlax International Journal of Education, 12(4), 49--59. \url{https://doi.org/10.34293/education.v12i4.7855}
\bibitem[Yasfi(2026)]{ref40} Yasfi, A. (2026). Future direction of artificial intelligence in higher education assessment: A systematic review. Cogent Education, 13, 2728183. \url{https://doi.org/10.1080/2331186X.2026.2728183}
\bibitem[Yeadon et al.(2026)]{ref12} Yeadon, W., Hardy, T., Mackay, P., \& Agra, E. (2026). Criterion-referenceability determines LLM-as-a-judge validity across physics assessment formats. arXiv:2603.14732 (current public record, accessed 2 November 2026; specific version label not independently confirmed). \url{https://arxiv.org/abs/2603.14732}
\bibitem[Zhang et al.(2026)]{ref54} Zhang, C., Ju, K. P., Lu, X., Yen, Y.-C. G., \& Rzeszotarski, J. M. (2026). EvaluAId: Human-AI Collaborative Evaluation of Open-Ended Student Essays. Proceedings of the CHI Conference on Human Factors in Computing Systems, Article 867, 1--20. \url{https://doi.org/10.1145/3772318.3790814}
\bibitem[Zhao et al.(2025)]{ref15} Zhao, C., Silva, M., \& Poulsen, S. (2025). Autograding mathematical induction proofs with natural language processing. International Journal of Artificial Intelligence in Education, 35, 3202--3232. \url{https://doi.org/10.1007/s40593-025-00498-2}
